\titledquestion{Balancing Act}

\textbf{Balancing Act}

A sequence of parentheses is \textit{balanced} if every open parenthesis is
matched by a later close parenthesis, and no close parenthesis appears before
its matching open. We represent such a sequence as a list:

\begin{codeblock}
  datatype paren = L | R

  type parens = paren list
\end{codeblock}

where \code{L} stands for \code{(} and \code{R} stands for \code{)}.

So, for instance:
\begin{itemize}
  \item \code{[L, R]} is balanced, representing \code{()}
  \item \code{[L, L, R, R]} is balanced, representing \code{(())}
  \item \code{[L, R, L, R]} is balanced, representing \code{()()}
  \item \code{[]} is balanced, representing the empty sequence
  \item \code{[R, L]} is \textbf{not} balanced
  \item \code{[L, L, R]} is \textbf{not} balanced
\end{itemize}

\begin{parts}
  \part[10]\hfill

    Implement \code{balanced}:

    \spec
      {balanced}
      {\code{parens -> bool}}
      {\code{true}}
      {\code{balanced ps} $\eeq$ \code{true} if and only if \code{ps} is
      balanced, and \code{false} otherwise}

    \begin{constraint}
      Your solution must traverse the list \textbf{exactly once}, and that
      traversal must be a single call to \code{foldl}. You may not use
      explicit recursion, \code{rev}, or any auxiliary list.
    \end{constraint}

    \textbf{Hint:} Your two counterexamples above tell you two different things
    that can go wrong. What does your accumulator need to remember in order to
    detect \textit{both} of them?

    \begin{solutionorbox}[26em]
      The accumulator must track the running depth \textit{and} whether the
      sequence has already gone irrecoverably wrong. \code{int option} does
      this: \code{SOME d} means ``still valid, at depth \code{d}'', and
      \code{NONE} means ``already broken, and nothing can fix it''.

      \begin{codeblock}
        fun balanced ps =
          case
            foldl
              (fn (_, NONE) => NONE
                | (L, SOME d) => SOME (d + 1)
                | (R, SOME d) =>
                    if d = 0 then NONE else SOME (d - 1))
              (SOME 0)
              ps
          of
            SOME 0 => true
          | _ => false
      \end{codeblock}

      Note that the post-processing step is essential: ending in
      \code{SOME d} for \code{d > 0} means unmatched open parens remain, as in
      \code{[L, L, R]}.
    \end{solutionorbox}

  \newpage

\end{parts}

\newpage
